\chapter{Curated EasyCrypt Proof-Object Case Studies}
\label{app:curated-easycrypt-case-studies}

This appendix replaces the least persuasive part of a purely generated
correspondence appendix with selected proof-object analyses.  The archived
exhaustive catalog remains in \ec{dissertation/etc/}; this appendix explains the
representative EasyCrypt objects that carry the main mathematical burden.

\section{Top-level theorem object}

The theorem family around \ec{euf\_cma\_security} and
\ec{haetae\_euf\_counted\_rom\_bound} is the formal version of the final
cryptographic claim:
\[
  \Adv_{\haetae,\rom}^{\eufcma}(A)
  \le
  \ec{haetae\_euf\_counted\_rom\_bound}.
\]
Mathematically, this is a universally quantified statement over adversary
modules satisfying the game interface and query-budget predicates.  Its proof is
not a single monolithic argument; it composes correctness, game-hop, sampler,
ROM, rejection, and reduction lemmas.

\section{Experiment module}

The signature ROM experiment in \ec{Sig\_ROM.ec} is read as a probabilistic
kernel:
\[
  \llbracket \mathsf{EUF\mbox{-}CMA}^{H,O}_{\haetae}(A)\rrbracket
  : S\to\mathcal{D}(S\times\{0,1\}).
\]
Its return bit is the forgery event.  This is the formal object whose return
probability defines the adversary's advantage.

\section{Lazy random oracle module}

The concrete lazy ROM in \ec{HAETAE\_ROM.ec} is a stateful oracle model.  Its
state contains the finite map of sampled answers and the logs needed by the ROM
programming proof.  The mathematical invariant is consistency:
\[
  q\in\operatorname{dom}(H_{\mathsf{tab}})
  \Rightarrow
  H(q)=H_{\mathsf{tab}}(q),
\]
and freshness:
\[
  q\notin\operatorname{dom}(H_{\mathsf{tab}})
  \Rightarrow
  H(q)\leftarrow D_H
\]
for the appropriate output distribution \(D_H\).

\section{ROM freshness lemma}

The sampler-prequery instrumentation proves a statement of the form
\[
  \Pr[\mathsf{sampler\_expand\_query}(\rho)\in Q_H]
  \le
  L_{\mathsf{ROM\mbox{-}pre}}.
\]
This is the exact mathematical content behind the informal phrase ``the simulator
programs a fresh random-oracle point.''  The important point is that \(Q_H\) is
the adversary's logged hash-query set, not an informal meta-level set.

\section{Sampler exactness lemma}

The transition from
\ec{ROSigningAttemptPaperSimSampler(HAETAE\_RO.FRO)} to
\ec{ROExactHyperballPaperSimSampler(HAETAE\_RO.FRO)} represents a statistical
comparison between an executable signing-side sampler and the exact paper
sampling distribution:
\[
  \Delta(\mathsf{Attempt},\mathsf{Exact})\le L_{\mathsf{sample}}.
\]
The proof obligations are support membership, rejection conditioning, and ROM
coverage.  This is where the proof pays the sampler and rejection losses.

\section{Active signing Hoare/equiv case}

The active signing oracle proof is centered on \ec{O.sign}.  The paper-level
statement is a clean-conditioned one-call theorem:
\[
  \Pr[O_{\mathsf{attempt}}.\mathsf{sign}:E\wedge C]
  \le
  \Pr[O_{\mathsf{exact}}.\mathsf{sign}:E]+\epsilon_{\mathsf{one}}.
\]
The condition \(C\) includes sampler coverage, public-key consistency, ROM
consistency, and absence of sampler prequery.  The checked proof matters because
chosen-message security requires this invariant to survive adaptive signing
queries.

\section{Reduction case}

The bimodal/MSIS reduction modules convert a successful final-game forgery into
a solver for the declared hardness relation.  Mathematically, the reduction is a
module transformer
\[
  A \mapsto B_A
\]
with a probability inequality
\[
  \Pr[G_{\mathsf{final}}(A):\mathsf{Forge}]
  \le
  \Adv_{\mathsf{Hard}}(B_A)+L_{\mathsf{extract}}.
\]
The corresponding EasyCrypt evidence appears in the reduction and hop-game
modules for bimodal self-target MSIS and Module-SIS lifting.

\section{Why these cases are enough for exposition}

The exhaustive declaration catalog is useful for audit, but a reader does not
learn the proof by reading thousands of templated semantic lines.  The case
studies above identify the proof objects that carry the cryptographic argument:
the final theorem, the experiment module, the lazy ROM, the sampler freshness
lemma, the sampler exactness lemma, the active signing proof, and the hardness
reduction.
